\chapter{The Number Splits Into Faces}

\section{One corridor, four faces}

Part II handed over a single object: a bracketed, invariant residue, floored and repeatable. This chapter asks what that one
object \emph{is}, and finds that the question has more than one answer --- not because the object is ambiguous but because it
presents a different face depending on where it is read. The residue is one corridor; charge, mass, phase, and curvature are
four faces of it; and which face you see is fixed by which handling of the tensor you perform. The object does not change.
The reading does.

The geometry of ``four faces'' is the decomposition of one tensor by the operations available to it. Take the residue as a
curvature and ask what invariants a single such object yields. Contract two indices and you get a trace, the Ricci part,
which the physical gauge reads as \emph{mass}-energy through $R_{\mu\nu} - \tfrac12 R g_{\mu\nu} = 8\pi T_{\mu\nu}$. Integrate
its dual flux through a closed surface and you get a conserved total, which the gauge reads as \emph{charge}, $Q = \oint
\star F$. Integrate its potential around a closed loop and you get a holonomy, which the gauge reads as \emph{phase},
$\exp\!\big(\tfrac{ie}{\hbar}\oint A\big)$. Leave it uncontracted, as the trace-free remainder, and it is \emph{curvature}
proper, the tidal shape --- the Weyl part, which no local source fixes. Four operations on one tensor: contract, flux, loop,
and leave-alone. Four faces, one corridor.

The unity underneath is that these are not four objects but four \emph{readings of index structure}. The residue is a tensor
with indices; every face is a decision about which indices to sum, which to integrate, and which to leave free. Mass is the
face where you trace; charge is the face where you flux; phase is the face where you loop; curvature is the face where you
touch nothing. This is why the instrument needs only one residue: the multiplicity of the physical world's names --- that
charge and mass and phase seem such different quantities --- is a multiplicity of \emph{handlings}, not of things. Noether's
theorem is the deep statement of this economy: each conserved face corresponds to a symmetry of the one action, and the
faces are its currents, not independent posits~\cite{noether1918}. One symmetry, one current, one face; one action beneath
them all.

It must be marked which of this is derived and which is read. That the residue is one invariant admitting several index
handlings --- that is the construction's, and exact. That a particular handling \emph{is} charge, or mass, or phase --- that
is the physical gauge's reading, laid over the handling and marked as such. The instrument does not prove that its trace is
the world's mass; it proves that its residue has a trace, and the physics recognizes mass in that trace. The book keeps the
distinction so that no face is ever mistaken for a derivation of the magnitude the world measures for it. The corridor is
built; the faces are read; and the reading is honest exactly because it is labeled a reading.

In the two verbs, the four faces are four tanges taken from one funge. The residue is the funge --- the world's covariant
invariant, bagged whole. To read it as charge, or mass, or phase, is to tange one characteristic out of that bag: to select
the trace, or the flux, or the loop, and leave the rest. The faces are not four bags; they are four selections from a single
bag. So the number splits into faces the way one tensor splits into its invariants --- by which index you choose to read ---
and the next section names the place where that choice is made, which is exactly where the four gauges of the whole work part
company.

\section{The partition point}

The last section showed one residue wearing four faces. This section names the place where the faces are assigned --- the
partition point, the site in the construction where the single corridor is cut into readings. It is a structural place, and
it is the same place where the four volumes of this work divide their labor. No single gauge owns the whole instrument; each
reads the corridor from its own side of the partition; and the physical gauge, this volume, owns exactly the faces where the
residue presents as force, orbit, field, and coupling. Naming the partition is naming why there must be four books and not
one.

The partition is, geometrically, the choice of projection. A tensor has no single canonical decomposition until you fix a
structure to decompose it against --- a frame, a splitting of the space into pieces, a set of projectors $P_i$ with $\sum_i
P_i = \mathbb{1}$. Each projector reads out one face: $P_{\text{trace}}$ gives the mass face, $P_{\text{flux}}$ the charge
face, and so on. The partition point is where these projectors are chosen. Before it, there is one residue; after it, there
are four readings; and the choice of projectors is precisely the choice of which face to make manifest. The completeness
relation $\sum_i P_i = \mathbb{1}$ is the statement that the four faces together exhaust the residue --- that nothing is lost
in the splitting, only distributed --- so the partition does not discard the corridor, it reads it out along four
complementary channels.

That no gauge owns the whole is a theorem about projections, not a modesty. A single projector $P_i$ sees only its own face
and is blind to the others; $P_{\text{trace}}$ cannot read the loop, $P_{\text{loop}}$ cannot read the trace. To recover the
whole residue you need all four projectors, and no one of them suffices. This is why the physical gauge, however complete on
its own faces, cannot claim the corridor entire: it reads force and field and coupling exactly, and it is structurally blind
to the faces the other gauges read, because those live in projections it does not perform. The instrument is genuinely
four-faced, and a reader who wanted the whole would need all four books, each reading one complementary channel of the one
residue.

There is a reason the partition is a strength rather than a fracture, and it is the reason the whole four-volume instrument
can be trusted at the end. Because the four projectors are complementary --- because they read the residue along channels
that do not overlap --- their readings are \emph{decorrelated}. What the mass face sees, the phase face cannot see; the four
readings share no vocabulary and no machinery. So when all four close, at the end, on the same bracket, the agreement cannot
be an artifact of a shared method, because there is no shared method: the projectors are orthogonal. The partition that
splits the corridor into four blind faces is exactly what makes their eventual agreement meaningful. A single face could be
tuned to bracket anything; four orthogonal faces cannot be tuned to agree by accident.

In the two verbs, the partition point is where one funge is divided among four tanges. The residue is funged whole --- one
covariant invariant --- and the partition assigns each gauge a tange, a projector, a characteristic to select. The physical
gauge tanges the force face, the field face, the coupling face; the other gauges tange theirs; and the projectors'
completeness, $\sum_i P_i = \mathbb{1}$, guarantees that the four tanges together recover the whole funge. So the corridor
partitions into faces the way a bagged invariant partitions among orthogonal selectors --- each gauge selecting its own
characteristic, none owning the bag alone. The next section makes exact the sense in which the face depends on where the
corridor is viewed, and why relabeling one face into another is the operation the following chapters will read as angular
momentum.

\section{Same residue, different reading}

The partition assigns the faces; this closing section examines the assignment itself and draws the lesson the whole part
exists for: the face is not in the residue but in the \emph{viewing}. The same invariant, read from a different frame,
presents as a different quantity, and the change from one face to another is a physical operation with a physical cost. This
is the seed of the next chapter's rotation and the one after's angular momentum: to relabel a face is to turn the corridor,
and turning the corridor is not free.

The mechanism is index raising and lowering, the metric trading one kind of reading for another. A single tensor component
reads one way with an index up and another with it down, and the metric converts between them, $V_\mu = g_{\mu\nu} V^\nu$.
The information is identical --- $V_\mu$ and $V^\nu$ encode the same invariant --- but the two are different readings, a
covariant one and a contravariant one, of one object. Push this through a change of frame and the same residue presents its
mass face to one observer and its phase face to another, exactly as the same electromagnetic field is a pure electric field
to one observer and a mixed electric-magnetic field to another boosted relative to the first. The field did not change; the
frame did; and the face followed the frame. What a reading \emph{is} depends on where the corridor is viewed.

This is worth stating as sharply as the geometry allows, because it is the opposite of the naive picture. Naively, mass and
charge and phase are different substances a body possesses, and measurement reads off how much of each it has. The
construction says instead that there is one residue, and mass, charge, and phase are the readings it yields under different
projections and frames --- that the ``substances'' are viewpoints. This does not make them unreal; the electric field an
observer reads is perfectly real to that observer. It makes them \emph{relational}: real, repeatable, and dependent on the
frame that reads them. The invariant is the residue; the faces are its frame-dependent readings; and the physics lives in
both --- in the invariant that all frames share and in the readings each frame takes.

The construction supplies the frame-dependence; the identification of specific frames with specific physical observers is the
gauge's reading, marked as such. What the instrument derives is that one invariant reads differently under different
projections and boosts. What the physics adds is that this is exactly how a Lorentz boost mixes electric and magnetic fields,
how a change of gauge shifts a phase, how a change of chart reassigns the coordinates of a curvature. The book keeps the two
apart: the frame-dependence of the reading is structural, and the recognition of it in the specific transformations of
electromagnetism is the physical gloss. Both are true; only one is derived; and the marking keeps the reader from mistaking
the gloss for a magnitude the instrument produced.

In the two verbs, ``same residue, different reading'' is one funge presenting under different tanges. The residue is the
single covariant bag; the frame is the tange, the contravariant choice of how to select from it; and changing the frame
changes which characteristic the tange pulls out, hence which face the funge presents. The invariant is what every tange
shares; the face is what each tange selects. So the corridor is read differently from different sides because the funge is
one and the tanges are many --- and because a tange is a choice, relabeling one face into another is an operation, with a
direction and a price. That operation is rotation, and the price is angular momentum; the next two chapters read them in
order, beginning with the named, stable face the corridor settles into when the counting forces it into one box: the
electron.
